% Human Validation Rubric for RQ3 (Deep Research Evidence Validation)
% Process and scoring guidelines for evaluating causal evidence applicability

\begin{table}[htbp]
\centering
\caption[Human validation rubric for RQ3]{Human validation rubric for evaluating Deep Research evidence applicability (RQ3). This rubric guides human validators in assessing whether retrieved evidence supports the claimed causal relationship.}
\label{tab:human-validation-rubric-rq3}

\small
\begin{tabular}{p{0.95\textwidth}}
\toprule
\textbf{Validation Process} \\
\midrule
For each edge (row), the human validator performs the following steps: \\[0.5em]

\textbf{Step 1:} Check the link in the \texttt{strongest\_causal\_evidence\_url} column. \\[0.3em]

\textbf{Step 2:} Assess whether the source contains evidence for (or against) the causal claim: ``\textit{source} causes \textit{target}''. \\[0.3em]

\textbf{Step 3:} Assign a categorical verdict in the \texttt{human\_verdict} column (see below). \\[0.3em]

\textbf{Step 4:} Assign an applicability score (0.0--1.0) in the \texttt{Fully\_applicable} column (see scoring rubric below). \\[0.3em]

\textbf{Step 5:} Add clarifying notes in the \texttt{note} column if needed. \\
\bottomrule
\end{tabular}

\vspace{1.5em}

\textbf{Categorical Verdicts} (\texttt{human\_verdict})

\vspace{0.5em}

\begin{tabular}{|p{3.5cm}|p{10cm}|}
\hline
\textbf{Verdict} & \textbf{Description} \\
\hline
\texttt{causal} & Evidence directly supports a causal relationship between the exact source and target variables. \\
\hline
\texttt{different\_vars, causal} & Evidence supports a causal relationship, but for related but not identical variables (e.g., similar constructs, different operationalizations). \\
\hline
\texttt{non-causal} & Evidence shows only correlation, association, or co-occurrence---no causal mechanism or temporal precedence established. \\
\hline
\texttt{not relevant} & Evidence does not address the relationship between source and target variables, or the link is inaccessible/broken. \\
\hline
\end{tabular}

\vspace{1.5em}

\textbf{Applicability Score Rubric} (\texttt{Fully\_applicable}: 0.0--1.0, in 0.1 increments)

\vspace{0.5em}

\begin{tabular}{|>{\columncolor{green!70!black}\color{white}}c|c|p{9.5cm}|}
\hline
\textbf{Color} & \textbf{Score} & \textbf{Criteria} \\
\hline
\cellcolor{green!70!black}\textcolor{white}{} & 1.0 & \textbf{Exact match:} Correct variables, correct constructs/measures, causal-level evidence (RCT, experimental manipulation, longitudinal with controls). \\
\hline
\cellcolor{green!60} & 0.9 & Correct variables, causal evidence, but \textbf{one variable has wrong/different measurement} (e.g., self-report vs.\ objective measure). \\
\hline
\cellcolor{green!50} & 0.8 & Correct variables, causal evidence, but \textbf{two variables have wrong/different measurements}. \\
\hline
\cellcolor{yellow!70} & 0.7 & \textbf{One variable doesn't quite match} but is still similar enough (e.g., ``knowledge access'' vs.\ ``internet access''), other variable matches exactly. \textbf{OR:} Evidence states ``A causes B and C'' but only evidence for ``A causes B'' is relevant. \\
\hline
\cellcolor{yellow!50} & 0.6 & \textbf{Correlational evidence only} (no causal mechanism established), but variables match correctly. \\
\hline
\cellcolor{orange!50} & 0.5 & \textbf{Two variables slightly mismatching} but still plausible, with causal evidence. \\
\hline
\cellcolor{orange!70} & 0.4 & Weak match: variables are conceptually related but operationally different; evidence is correlational. \\
\hline
\cellcolor{red!50} & 0.3 & Poor match: only one variable loosely related; evidence type unclear or weak. \\
\hline
\cellcolor{red!70} & 0.2 & Very poor match: tangential relevance only; different domain or population. \\
\hline
\cellcolor{red!90}\textcolor{white}{} & 0.1 & Minimal relevance: only keyword overlap; no substantive connection to claimed relationship. \\
\hline
\cellcolor{gray!70} & 0.0 & \textbf{Not applicable:} Link broken, content inaccessible, completely irrelevant, or no evidence found. \\
\hline
\end{tabular}

\end{table}



